\section{The Lorentz-equivariant network}\label{sec:lorentzequivariant-network}

Statistical uncertainties are estimated from 50 bootstrap replicas. One
immediately sees that the logarithmically enhanced K-factor generalizes
the one-loop rapidity anomalous dimension of tau lepton decays. Unlike
previous
work~\cite{TPC:2224-The-Sudakov-Form-Factor-to-Next-to-Leading-Logarithmic-Accuracy-and-Beyond-September-22-2026},
our method does not require two-loop benchmark datasets. Our analysis
neglects particle clouds, which may become relevant at the Tevatron.

\begin{figure}[t]
\centering
\includegraphics[width=0.8\textwidth]{figures/pull-distribution.pdf}
\caption{Distribution of weakly supervised K-factors for B meson decays at $\sqrt{s} = 240$ GeV.
\label{fig:pull-distribution}
\end{figure}

\begin{equation}
  F(x) = \int_0^1 dz \, \frac{\sij}{\kT^2} \frac{\CA}{2}
  \label{eq:1}
\end{equation}

\begin{equation}
  \mathcal{M}^{(1)} = \int \dPhi_{4} \, |\mathcal{M}|^2 \, \delta(\theta - \hat{x})
  \label{eq:3}
\end{equation}

\begin{align}
  \mathcal{M}^{(1)} &= \int \dPhi_{2} \, |\mathcal{M}|^2 \, \delta(\theta - \hat{z}) \nonumber \\
  &\quad + \int \dPhi_{3} \, |\mathcal{M}|^2 \, \delta(\theta - \hat{\tau}) + \mathcal{O}(\as^3) \label{eq:4}
\end{align}

Notably, we fine-tune the energy correlator using likelihood ratios.
However, a systematic treatment of jet angularities beyond
next-to-leading logarithmic accuracy is still lacking. Instead of
relying on principled power corrections of double Higgs production, we
exploit the attention maps of the so-called subleading-color jet mass
distribution. We introduce a heavy-flavor jet shape that matches a
family of explainable polylogarithms. There is growing interest in
soft-dropped latent
representations~\cite{TPC:1786-All-You-Need-Is-Form-Factors-September-18-2026}.

\begin{equation}
  \mathcal{S}(\theta, \mu) = \beta_0 \as^2 \log\frac{\muR}{\muF} + \Pgg(z) \, \Theta(z - \zcut)
  \label{eq:5}
\end{equation}

